\begin{table}
\caption{Tabela 1. As duas bacias: código de Otto Pfafstetter (BHB250, nível 3), área, células ERA5 que as intersectam e que entram na comparação final, faixa de latitudes dos centros das células, participação dos biomas (IBGE) e de Terras Indígenas ou Unidades de Conservação federais, sem contar duas vezes as áreas sobrepostas.}
\label{tab:02_regioes::comparacao}
\begin{tabular}{llrrrlrrr}
\toprule
Bacia & Código Otto & Área (mil km²) & Células & Células na comparação & Latitudes & Amazônia (\%) & Cerrado (\%) & TIs ou UCs (\%) \\
\midrule
Xingu & 642 & 509 & 779 & 668 & -15,00° a -1,50° & 93 & 7 & 53 \\
Tapajós & 644 & 493 & 762 & 652 & -15,00° a -2,25° & 80 & 20 & 35 \\
\bottomrule
\end{tabular}
\end{table}
